\sectionnonum{Rezultatai}
Darbo metu buvo sukurtas RISC-V architektūros emuliatorius, palaikantis Linux operacinės sistemos paleidimą, pagrindinius privilegijuoto režimo mechanizmus, pertrauktis, virtualius įrenginius bei derinimą naudojant GDB nuotolinio derinimo protokolą. Emuliatoriaus veikimas buvo tikrinamas naudojant \texttt{riscv-tests}, savarankiškai sukurtus ELF testus bei Linux sistemos paleidimą. Realizuotos papildomos diagnostinės instrukcijos skirtos emuliatoriaus būsenos palyginimui su Spike realizacija.

Darbo metu taip pat išplėsta RISC-V instrukcijų architektūra pridedant QPU instrukcijas kvantinių skaičiavimų vykdymui. Sukurtas instrukcijų dekodavimo ir vykdymo mechanizmas leidžia svečio sistemoje naudoti \texttt{libquantum} biblioteką. Atlikti našumo matavimai parodė, kad vykdant skaičiavimus su didesniu kubitų skaičiumi emuliacijos suletėjimas tampa santykinai mažas, nes didžiąją vykdymo laiko dalį sudaro pačios kvantinės bibliotekos atliekami skaičiavimai.

Instrukcijų talpyklos našumo tyrimas parodė, kad pasirinkta talpyklos realizacija turi didelę įtaką bendram emuliatoriaus našumui. Paprasta masyvu grįsta talpykla šiam emuliatoriui yra efektyvesnė už maišos lentele grįstą realizaciją. Gauti rezultatai parodo, kad net ir nesudėtingi instrukcijų spartinimo metodai gali reikšmingai sumažinti Linux sistemos užkrovimo laiką bei pagerinti bendrą emuliatoriaus veikimo efektyvumą.

\sectionnonum{Išvados}

Šiame darbe sukurtas emuliatorius parodė, kad pilnai veikiančiai RISC-V programinės įrangos grandinei nebūtina modeliuoti visos realaus procesoriaus mikroarchitektūros. Darbe pasirinkta siauresnė, bet praktiška kryptis: realizuoti tuos procesoriaus, privilegijų, virtualios atminties, pertraukimų ir įrenginių mechanizmus, kurių Linux branduoliui vykdyti. Toks požiūris leido susieti atskiras architektūros dalis į veikiančią sistemą ir kartu išlaikyti emuliatoriaus realizaciją pakankamai aiškią bei kontroliuojamą.

Eksperimentinis QPU plėtinys parodė, kad RISC-V architektūrą galima naudoti ne tik standartinių instrukcijų emuliavimui, bet ir specializuotų skaičiavimo modelių tyrimui. QPU instrukcijos buvo integruotos į bendrą instrukcijų dekodavimo ir vykdymo kelią, o jų būsena prijungta prie Linux procesų valdymo. Tai leido kvantinių skaičiavimų simuliavimą pateikti kaip svečio sistemos instrukcijų rinkinio dalį, o ne tik kaip atskirą šeimininko sistemoje veikiančią biblioteką.

Sąmoningai ribota emuliatoriaus apimtis gali būti naudinga ir mokymo tikslais. Lyginant su didesniais, universalesniais emuliatoriais, tokia sistema yra lengviau skaitoma ir leidžia studentams tiesiogiai matyti ryšį tarp įvairių sistemos dalių: instrukcijų vykdymo, CSR, išimčių, adresų vertimo, MMIO ir operacinės sistemos paleidimo. Tokį emuliatorių būtų galima naudoti įvairaus sudėtingumo užduotims: naujų instrukcijų projektavimas ir implementacija, naujo MMIO įrenginio kūrimas, RISC-V programų rašymas \textit{bare-metal} aplinkoje. Taip pat galima tirti arba derinti Linux paleidimo eigą: pateikus nepilnai įgyvendintą emuliatorių prašoma implementuoti trūkstamas dalis, jas identifikuojant GDB derinimo pagalba. Dėl to darbas gali būti laikomas ne tik konkretaus emuliatoriaus realizacija, bet ir pagrindu tolimesniems architektūros, operacinių sistemų bei specializuotų instrukcijų plėtinių eksperimentams ir mokymams.
